\begin{pastexam}[label={px:nsysu-108-dmalg-04}]{中山資工所資安組 108 學年度 第 8 題}
（計算）

(Algorithm points) For each of the following pairs $a,b\in\mathbf{Z}^+$, determine $\gcd(a,b)$ (by
Euclidean Algorithm) and express it as a linear combination of $a,b$.
\begin{enumerate}[label=(\alph*)]
\item {[6\%]} 251, 1920
\item {[6\%]} 1371, 2587
\item {[6\%]} 1689, 6001
\end{enumerate}
\end{pastexam}

\begin{pastsoln}
答案：三組的最大公因數都是 $1$。
\begin{itemize}
\item (a) $1 = (-589) \cdot 251 + 77 \cdot 1920$；
\item (b) $1 = (-1185) \cdot 1371 + 628 \cdot 2587$；
\item (c) $1 = (-2615) \cdot 1689 + 736 \cdot 6001$。
\end{itemize}

\textbf{做法}。 輾轉相除法（資結書虛擬碼~6.9、6.10）：$\gcd(x, y) = \gcd(y, x \bmod y)$，一直做到餘數是 $0$，
最後一個不是 $0$ 的餘數就是最大公因數。 要寫成線性組合，在相除的同時記下每個餘數是兩個原數的幾倍相加
（擴展歐幾里得算法）：令大的數是 $B$、小的數是 $A$，
\[
r_0 = B,\quad r_1 = A,\quad r_{i+1} = r_{i-1} - q_i r_i \quad (q_i = \lfloor r_{i-1}/r_i \rfloor),
\]
每一列都維持 $r_i = x_i B + y_i A$。 開頭 $(x_0, y_0) = (1, 0)$、$(x_1, y_1) = (0, 1)$，之後 $x$、$y$ 用和 $r$
一樣的式子更新：$x_{i+1} = x_{i-1} - q_i x_i$，$y_{i+1} = y_{i-1} - q_i y_i$。 因為三個式子是同一個線性運算，
等式 $r_i = x_i B + y_i A$ 一路保持。 做到 $r = 1$ 那一列，就是答案。

\textbf{(a)} $B = 1920$，$A = 251$：
\begin{center}
\begin{tabular}{rrrrr}
\toprule
$i$ & $r_i$ & $q_i$ & $x_i$ & $y_i$ \\
\midrule
$0$ & $1920$ & & $1$ & $0$ \\
$1$ & $251$ & $7$ & $0$ & $1$ \\
$2$ & $163$ & $1$ & $1$ & $-7$ \\
$3$ & $88$ & $1$ & $-1$ & $8$ \\
$4$ & $75$ & $1$ & $2$ & $-15$ \\
$5$ & $13$ & $5$ & $-3$ & $23$ \\
$6$ & $10$ & $1$ & $17$ & $-130$ \\
$7$ & $3$ & $3$ & $-20$ & $153$ \\
$8$ & $1$ & $3$ & $77$ & $-589$ \\
\bottomrule
\end{tabular}
\end{center}
相除的過程：$1920 = 7 \cdot 251 + 163$，$251 = 1 \cdot 163 + 88$，$163 = 1 \cdot 88 + 75$，$88 = 1 \cdot 75 + 13$，
$75 = 5 \cdot 13 + 10$，$13 = 1 \cdot 10 + 3$，$10 = 3 \cdot 3 + 1$，$3 = 3 \cdot 1 + 0$。 最後一個非零餘數是 $1$，
$\gcd(251, 1920) = 1$。 驗算：$77 \cdot 1920 = 147840$，$589 \cdot 251 = 147839$，相減得 $1$。

\textbf{(b)} $B = 2587$，$A = 1371$：
\begin{center}
\begin{tabular}{rrrrr}
\toprule
$i$ & $r_i$ & $q_i$ & $x_i$ & $y_i$ \\
\midrule
$0$ & $2587$ & & $1$ & $0$ \\
$1$ & $1371$ & $1$ & $0$ & $1$ \\
$2$ & $1216$ & $1$ & $1$ & $-1$ \\
$3$ & $155$ & $7$ & $-1$ & $2$ \\
$4$ & $131$ & $1$ & $8$ & $-15$ \\
$5$ & $24$ & $5$ & $-9$ & $17$ \\
$6$ & $11$ & $2$ & $53$ & $-100$ \\
$7$ & $2$ & $5$ & $-115$ & $217$ \\
$8$ & $1$ & $2$ & $628$ & $-1185$ \\
\bottomrule
\end{tabular}
\end{center}
相除：$2587 = 1 \cdot 1371 + 1216$，$1371 = 1 \cdot 1216 + 155$，$1216 = 7 \cdot 155 + 131$，
$155 = 1 \cdot 131 + 24$，$131 = 5 \cdot 24 + 11$，$24 = 2 \cdot 11 + 2$，$11 = 5 \cdot 2 + 1$，$2 = 2 \cdot 1 + 0$。
$\gcd = 1$。 驗算：$628 \cdot 2587 = 1624636$，$1185 \cdot 1371 = 1624635$，相減得 $1$。

\textbf{(c)} $B = 6001$，$A = 1689$：
\begin{center}
\begin{tabular}{rrrrr}
\toprule
$i$ & $r_i$ & $q_i$ & $x_i$ & $y_i$ \\
\midrule
$0$ & $6001$ & & $1$ & $0$ \\
$1$ & $1689$ & $3$ & $0$ & $1$ \\
$2$ & $934$ & $1$ & $1$ & $-3$ \\
$3$ & $755$ & $1$ & $-1$ & $4$ \\
$4$ & $179$ & $4$ & $2$ & $-7$ \\
$5$ & $39$ & $4$ & $-9$ & $32$ \\
$6$ & $23$ & $1$ & $38$ & $-135$ \\
$7$ & $16$ & $1$ & $-47$ & $167$ \\
$8$ & $7$ & $2$ & $85$ & $-302$ \\
$9$ & $2$ & $3$ & $-217$ & $771$ \\
$10$ & $1$ & $2$ & $736$ & $-2615$ \\
\bottomrule
\end{tabular}
\end{center}
相除：$6001 = 3 \cdot 1689 + 934$，$1689 = 1 \cdot 934 + 755$，$934 = 1 \cdot 755 + 179$，$755 = 4 \cdot 179 + 39$，
$179 = 4 \cdot 39 + 23$，$39 = 1 \cdot 23 + 16$，$23 = 1 \cdot 16 + 7$，$16 = 2 \cdot 7 + 2$，$7 = 3 \cdot 2 + 1$，
$2 = 2 \cdot 1 + 0$。 $\gcd = 1$。 驗算：$736 \cdot 6001 = 4416736$，$2615 \cdot 1689 = 4416735$，相減得 $1$。

線性組合不唯一：$x B + y A = 1$ 的一組解加上 $(-kA, kB)$ 仍是解，例如 (a) 也可以寫成
$1 = (-589 + 1920) \cdot 251 + (77 - 251) \cdot 1920 = 1331 \cdot 251 - 174 \cdot 1920$。 考卷寫出任何一組並驗算即可。
不想記表格的話，也可以從倒數第二個等式往回代（$1 = 10 - 3 \cdot 3$，再把 $3$、$10$、$13 \ldots$ 依序用前一個等式
換掉），結果相同，只是比較容易抄錯。
\end{pastsoln}
